\section{第10章：波动（10.5练习题）}
\sourceproblem{选择1}{179}{185}
\problemtext 比较静止波源的周期、质点振动速度及传播方向上质点相位关系；原选项将相位差限制为不超过 $\pi$。
\workings 对稳定单频右行波，$y=A\cos(\omega t-kx+\phi)$，周期等于波源周期，质点振动速度通常不等于波速。未折返的相位差为 $-kx$，下游振动有传播时延 $x/u$。但将相位差约化到 $(-\pi,\pi]$ 后，符号并非总为负，例如 $x=3\lambda/4$ 时差为 $-3\pi/2$，等价于 $+\pi/2$。
\crosscheck 按题中相位差约定，四个选项无普适正确者；若改为不折返的传播时延相位，教材意图为C。
\sourceproblem{选择2}{179}{185}
\problemtext $y=0.03\cos[6\pi(t+0.01x)]$（SI），辨认波的参数与方向。
\workings $A=0.03$ m，$\omega=6\pi$ rad/s，$T=1/3$ s，$k=0.06\pi$ rad/m，$u=\omega/k=100$ m/s，沿负x传播，选B。
\crosscheck 固定相位给 $dx/dt=-100$ m/s；注意0.01乘在x上。
\sourceproblem{选择3}{179}{185}
\problemtext $y=0.1\cos(3\pi t-\pi x+\pi)$（SI），图中a为上升过零点、b为相邻正峰，判断参数及相差。
\workings 振幅0.1 m，$\lambda=2$ m，$u=3$ m/s。a、b相隔四分之一波长，绝对相位差 $\pi/2$，选C。
\crosscheck O点初位移为负并不意味着振幅为负；传播方向为正x。
\sourceproblem{选择4}{179}{185}
\problemtext 图示A在正位移下降段、B在负位移上升段，已知A质元动能正在增加，判断行波方向和能量变化（图在PDF186）。
\workings 简谐质元 $a=-\omega^2y$，$\dot K=mva$。A处 $a<0$ 且 $\dot K>0$，故 $v<0$。因空间斜率也为负，满足左行波 $v=u\partial_x y$，选B。B处速度、加速度均为正，动能也增加。
\crosscheck 理想无耗散弹性行波的局部弹性势能密度与动能密度同相变化；不能将单质点弹簧振子的能量交换关系直接用于介质微元。
\sourceproblem{选择5}{180}{186}
\problemtext 负x方向传播的横波，按图判断A、B、C、D质元瞬时速度。
\workings 左行波满足 $v=u\partial_x y$；图中A、B、D位于下降段，速度负，C在上升段，速度正，故选D。
\crosscheck B虽在平衡位置，速率却为最大；位移为零不等于静止。
\sourceproblem{选择6}{180}{186}
\problemtext 频率500 Hz、波速360 m/s的波，求相位差 $\pi/3$ 对应的距离。
\workings $\lambda=u/\nu=0.72$ m，取最小正间距 $\Delta x=\lambda/6=0.12$ m，选D。
\crosscheck 若仅给模 $2\pi$ 的无向相位差，距离还可为 $|j\lambda\pm\lambda/6|$，j为整数；本选项按最小正间距。
\sourceproblem{选择7}{180}{186}
\problemtext 左行波初始波形在O点正斜率过零，求余弦式初相。
\workings O处 $y=0$ 且 $v=u\partial_x y>0$，故 $\cos\phi=0,\sin\phi<0$，$\phi=3\pi/2$（模 $2\pi$），选D。
\crosscheck 波形斜率与质元速度同号只适用于左行波。
\sourceproblem{选择8}{180}{186}
\problemtext 左行横波图中1、2、3分别为下降过零、波谷、上升过零，求再过四分之一周期时三点位移。
\workings 波形左移 $\lambda/4$，固定点取原波形右侧四分之一波长处的值；依次为 $-A,0,+A$，选B。
\crosscheck 质元横向振动，不随波形整体向左平移。
\sourceproblem{选择9}{180}{186}
\problemtext 图示右行波在O点正斜率过零，求余弦初相。
\workings $v=-u\partial_x y<0$，且 $y=0$，故初相 $\pi/2$（模 $2\pi$），选B。
\crosscheck 与选择7方向相反，初速度及初相不同。
\sourceproblem{选择10}{180}{186}
\problemtext 波频100 Hz、波速300 m/s，两点相距小于波长且相位差 $\pi/3$，求距离。
\workings $\lambda=3$ m。若相位差指未折返的绝对差，则 $\Delta x=0.5$ m，选C。若仅指模 $2\pi$ 的最小无向相差，则 $2.5$ m亦满足。
\crosscheck 原选项仅含0.5 m；小于一波长并不足以排除第二种模相位解，须说明相位约定。
\sourceproblem{选择11}{181}{187}
\problemtext 频率4 Hz的右行波，两点开始振动相隔0.25 s，求相位差。
\workings 传播时延对应未折返相差 $2\pi\nu\Delta t=2\pi$，选D。
\crosscheck 模 $2\pi$ 后两点同相；相差2π与同相并不矛盾。
\sourceproblem{选择12}{181}{187}
\problemtext 初始波形幅0.02 m、波长2 m，P位于x=1 m的下降过零点；P振动图周期0.2 s、初时负向过零，确定波方程。
\workings P点空间斜率和瞬时速度同为负，故波向负x传播。$\omega=10\pi$ rad/s，$k=\pi$ rad/m，$u=10$ m/s。O点初始空间斜率正、位移零，取 $\phi=-\pi/2$。于是 $y=0.02\cos[10\pi(t+x/10)-\pi/2]$ m，选B。
\crosscheck 代入 $x=1$ 得 $y_P=-0.02\sin10\pi t$，代入 $t=0$ 得 $0.02\sin\pi x$，同时符合两图。
\sourceproblem{选择13}{181}{187}
\problemtext 右行波初始图中P为下降过零点，周期2 s，选P点速度曲线。
\workings $v_P(0)=-u\partial_x y>0$，且过零时速率最大。故 $v_P(t)=\omega A\cos\omega t$，$\omega=\pi$ rad/s，选A。
\crosscheck 1 s速度为负极值，2 s回到正极值；图示量为振动速度。
\sourceproblem{选择14}{181}{187}
\problemtext 左行波速u，$x=b$ 处振动为 $A\cos(\omega t+\phi_0)$，求全波方程。
\workings 左行波同相变量为 $t+x/u$，用参考点b校准得 $y=A\cos\{\omega[t+(x-b)/u]+\phi_0\}$，选C。
\crosscheck $x=b$ 时恢复已知振动式；固定相位沿负x移动。
\sourceproblem{选择15}{181}{187}
\problemtext 波以2 m/s沿正x传播，原点振动 $0.06\cos(\pi t/5)$ m，求5 s的波形。
\workings $y=0.06\cos[(\pi/5)(t-x/2)]$，故 $t=5$ 时为 $y=0.06\cos[\pi(1-0.1x)]$ m，选B。
\crosscheck 波长20 m；仅在波已到达的区域适用稳态波函数。
\sourceproblem{选择16}{182}{188}
\problemtext $y=A\cos[2\pi(\nu t-x/\lambda)]$，在 $t=1/\nu$ 时比较 $x_1=3\lambda/4,x_2=\lambda/4$ 的质元速度。
\workings $v=-2\pi\nu A\sin[2\pi(\nu t-x/\lambda)]$。两处相位分别为 $\pi/2,3\pi/2$，速度分别为 $-2\pi\nu A,+2\pi\nu A$，故比为-1，选A。
\crosscheck 两速率相同但速度方向相反。
\sourceproblem{选择17}{182}{188}
\problemtext 给定 $y=4\cos(5\pi t+3\pi x)$，选择能与其形成驻波的波。
\workings 同方向偏振、同频、同幅且反向传播的波为 $y=4\cos(5\pi t-3\pi x)$，选C；合成 $8\cos3\pi x\cos5\pi t$。
\crosscheck 原波左行，所选波右行；沿不同坐标方向传播的波不构成该一维驻波。
\sourceproblem{选择18}{182}{188}
\problemtext 同一介质中两相干平面简谐波强度比4，求振幅比。
\workings 对同种线性波模式、相同频率及传播速度，平均强度 $I\propto A^2$，故 $A_1/A_2=2$，选C。
\crosscheck 相干波在重叠区域的总强度还含干涉项；题目比较两分波各自强度。
\sourceproblem{选择19}{182}{188}
\problemtext 两平面波源分别作 $A_1\cos\omega t$、$A_2\sin\omega t$ 振动，求等距点P的合幅。
\workings 等传播距离使初相差保持为 $\pi/2$，所以 $A_P=\sqrt{A_1^2+A_2^2}$，选C。
\crosscheck 取教材无衰减平面波假设，且两波在P处偏振方向相同。
\sourceproblem{选择20}{182}{188}
\problemtext 相干波到P的路程分别为 $2\lambda,2.2\lambda$，第一波源振动 $A\cos(2\pi t+\pi/2)$，P处相消干涉，求第二波源振动。
\workings 到达相差 $\phi_2-\pi/2-0.4\pi$，相消相位条件为奇数倍π，得 $\phi_2=1.9\pi\equiv-0.1\pi$。仅D的相位符合，但D振幅为2A，若两分波传播不衰减，P处最小合幅为A而非零。
\crosscheck 若相消指相位反相、干涉减弱，选D；若要求完全消光且分波幅值保持，则第二源应为 $A\cos(2\pi t-0.1\pi)$，原选项无严格正确者。若存在幅值衰减，还需其传播规律。
\sourceproblem{选择21}{182}{188}
\problemtext 两平面相干波源同幅2 cm且相差π，求中垂线上合振幅。
\workings 两程相同，到达相差仍为π，等幅反相抵消，合幅0，选A。
\crosscheck 若两波偏振不一致或到达幅值不等，则不能完全抵消。
\sourceproblem{选择22}{182}{188}
\problemtext 驻波相邻波节之间的质元振幅和相位关系。
\workings 驻波 $y=f(x)\cos(\omega t+\phi)$，相邻节点之间 $f(x)$ 同号，非节点质元同相而振幅随位置变化，选B。
\crosscheck 对称位置可能振幅相等；选项的振幅不同表示并非处处相同，而非任意两点必不同。
\sourceproblem{选择23}{182}{188}
\problemtext 两反向行波 $A\cos[2\pi(\nu t\mp x/\lambda)]$，求合成驻波振幅（选项续PDF189）。
\workings 和差化积得 $y=2A\cos(2\pi x/\lambda)\cos(2\pi\nu t)$，振幅取非负值 $|2A\cos(2\pi x/\lambda)|$，选D。
\crosscheck 负的空间系数表示相位反转，不表示负振幅。
\sourceproblem{选择24}{183}{189}
\problemtext 两波 $6\cos[\pi(0.02x-8t)/2]$ 和 $6\cos[\pi(0.02x+8t)/2]$（SI），求驻波节点。
\workings 合波 $y=12\cos(0.01\pi x)\cos4\pi t$，节点满足 $0.01\pi x=(j+1/2)\pi$，即 $x=50+100j$ m，选A。
\crosscheck 节点间距100 m，分波波长200 m；x取正负整数序列。
\sourceproblem{选择25}{183}{189}
\problemtext 汽笛750 Hz，机车以90 km/h远离静止观察者，声速340 m/s，求接收频率。
\workings 源速25 m/s，静止介质中远离时 $f'=f u/(u+v_s)=750(340/365)\approx698.63$ Hz，选B。
\crosscheck 远离时频率降低；移动声源与移动观察者的经典公式不同。
\sourceproblem{填空1}{183}{189}
\problemtext 右行横波，图示A处正斜率、B及C处负斜率，求三点瞬时运动方向。
\workings $v=-u\partial_x y$，故A向下，B、C向上。
\crosscheck C位于平衡点附近下降段，仍有向上速度。
\sourceproblem{填空2}{183}{189}
\problemtext 周期0.04 s、波速300 m/s的右行波，求x=10、16 m两质点相差。
\workings $\lambda=uT=12$ m，$\phi_{16}-\phi_{10}=-2\pi(6/12)=-\pi$，无向相差π。
\crosscheck 后一点在未折返相位中滞后半周期。
\sourceproblem{填空3}{183}{189}
\problemtext 右行波初始图中原点正峰、x=0.4 m处负峰，幅值0.2 m，波速100 m/s，求波长、振幅与频率。
\workings $\lambda/2=0.4$ m，故 $\lambda=0.8$ m，$A=0.2$ m，$\nu=u/\lambda=125$ Hz。
\crosscheck 相邻过零点0.2、0.6 m相隔半波长；正峰在0及约0.8 m。
\sourceproblem{填空4}{183}{189}
\problemtext 频率4 Hz的波，两点开始振动时间差0.25 s，求相差。
\workings 未折返相差大小 $2\pi\nu\Delta t=2\pi$；两点同相。
\crosscheck 有向相差取决于a、b在传播方向上的先后关系。
\sourceproblem{填空5}{183}{189}
\problemtext 在 $y=A\cos(\omega t-\omega x/u)$ 中解释 $x/u,\omega x/u,y$。
\workings 对 $x\ge0$ 的右行波，$x/u$ 为相对原点的传播时延；$\omega x/u$ 为该时延产生的未折返相位滞后量；y为位置x处质元在t时刻相对平衡位置的有向位移。
\crosscheck 对负x，$x/u$ 为带符号的相对时延；y不是波传播距离。
\sourceproblem{填空6}{183}{189}
\problemtext $y=0.025\cos(125t-0.37x)$（SI），求角频率、波速和波长。
\workings $\omega=125$ rad/s，$k=0.37$ rad/m，$u=125/0.37\approx337.84$ m/s，$\lambda=2\pi/0.37\approx16.98$ m。
\crosscheck 相位固定时沿正x传播，$u=\nu\lambda$。
\sourceproblem{填空7}{183}{189}
\problemtext 声纳波 $y=1.2\times10^{-3}\cos(3.14\times10^5t-220x)$（SI），求频率、波长、波速。
\workings $\nu=3.14\times10^5/(2\pi)\approx5.00\times10^4$ Hz，$\lambda=2\pi/220\approx0.02856$ m，$u=3.14\times10^5/220\approx1427.27$ m/s。
\crosscheck 三个量满足 $u=\nu\lambda$；频率属于超声范围。
\sourceproblem{填空8}{183}{189}
\problemtext 右行波振幅0.08 m、频率50 Hz、波长4 m，原点初时在正峰，求原点振动及波函数（题尾续PDF190）。
\workings $\omega=100\pi$ rad/s，$k=\pi/2$ rad/m。$y_0=0.08\cos100\pi t$ m，$y=0.08\cos(100\pi t-\pi x/2)$ m。
\crosscheck 波速200 m/s，代入x=0恢复原点振动。

\sourceproblem{填空9}{184}{190}
\problemtext 波长1 m、振幅0.1 m、周期0.5 s，右行波原点初相为零，求波函数。
\workings $\omega=4\pi$ rad/s，$k=2\pi$ rad/m，故 $y=0.1\cos(4\pi t-2\pi x)$（SI）。
\crosscheck 波速2 m/s，固定相位沿正x移动。
\sourceproblem{填空10}{184}{190}
\problemtext 右行波周期0.5 s、波长10 m、振幅0.1 m，源初时为正峰，求半波长处振动式及半周期时四分之一波长处速度。
\workings $y=0.1\cos(4\pi t-\pi x/5)$，故 $y(5,t)=-0.1\cos4\pi t$ m。$v=-0.4\pi\sin(4\pi t-\pi x/5)$，在 $t=0.25$ s、$x=2.5$ m 得 $v=-0.4\pi$ m/s。
\crosscheck 目标相位为 $\pi/2$；位移零、速度负且达最大速率。
\sourceproblem{填空11}{184}{190}
\problemtext 给定 $y=A\cos(Bt-Cx)$，三个常量均正，求波长、波速和相距d两点的相差。
\workings $\lambda=2\pi/C$，$u=B/C$；沿传播方向后一点相对前一点的未折返相差为 $-Cd$，滞后量为 $Cd$。
\crosscheck 相位可模 $2\pi$ 化简；B、C分别具有角频率、波数的量纲。
\sourceproblem{填空12}{184}{190}
\problemtext 图示 $t=T/4$ 时右行波，波速330 m/s，原点下降过零，x=1 m处负峰，x=3 m处正峰，振幅1.10 m，求波函数。
\workings $\lambda=4$ m，$k=\pi/2$ rad/m，$\omega=165\pi$ rad/s。此时波形 $y(x,T/4)=-1.10\sin(\pi x/2)$，故 $y=1.10\cos(165\pi t-\pi x/2+\pi)$ m。
\crosscheck 在 $t=T/4$、$x=1$ m 代入得到-1.10 m；原点此时速度为正，与右行波负空间斜率一致。
\sourceproblem{填空13}{184}{190}
\problemtext S点振动 $A\cos(\omega t+\phi)$，图a坐标正向与波速同向、图b坐标正向与波速反向，两图S坐标均为 $x_0$，求波函数。
\workings 图a为右行波，$y_a=A\cos\{\omega[t-(x-x_0)/u]+\phi\}$；图b为负x方向行波，$y_b=A\cos\{\omega[t+(x-x_0)/u]+\phi\}$。
\crosscheck 两式在 $x=x_0$ 均恢复已知振动；图b的x轴正方向指向图面左方，不能按纸面方向直接判号。
\sourceproblem{填空14}{184}{190}
\problemtext 同心球面半径 $R_1,R_2$，分别取面积 $\Delta S_1,\Delta S_2$，求平均能流之比。
\workings 若为均匀无吸收球面波，平均强度 $I=P/(4\pi R^2)$，故 $\overline P_1/\overline P_2=R_2^2\Delta S_1/(R_1^2\Delta S_2)$。
\crosscheck 若波源有方向性、介质吸收或面积跨不同辐射方向，则需相应强度积分；原题按均匀球面波模型。
\sourceproblem{填空15}{184}{190}
\problemtext 波源功率4 W，均匀无吸收球面波，求距源1 m的能流密度。
\workings 若输入功率全部转化为波的平均辐射功率，$I=4/(4\pi\cdot1^2)=1/\pi$ W每平方米。
\crosscheck 若波源另有内部耗散，应以实际辐射功率代替输入功率。
\sourceproblem{填空16}{184}{190}
\problemtext 均匀球面波发射功率4 W，距源2 m，原题问能源密度。
\workings 平均能流密度为 $I=4/[4\pi(2)^2]=1/(4\pi)$ W每平方米。若问体积平均能量密度，则对无耗散非色散行波 $\overline w=I/u=1/(4\pi u)$ J每立方米，需另知能量传播速度u。
\crosscheck 原页用词为能源密度，不能将W每平方米直接写为J每立方米；给定数据仅唯一确定能流密度。
\sourceproblem{填空17}{184}{190}
\problemtext 平面简谐机械行波中质元在t时刻总机械能10 J，求一个周期后振动动能。
\workings 对理想线性无耗散弹性行波，局部动能密度与弹性势能密度相等（例如弦波 $w_K=\rho_l y_t^2/2,\ w_U=F y_x^2/2$，且 $F=\rho_l u^2,\ y_t=-uy_x$）。故当时动能5 J；周期后各场量复现，动能仍5 J。
\crosscheck 按教材质元微小体积模型；有限区域积分在同一理想模型下亦成立。若是一般驻波或单独简谐振子，总能量不能单独确定瞬时动能。
\sourceproblem{填空18}{184}{190}
\problemtext 同相等幅波源至P距离3及4.5个波长，传播幅值不变，求合幅。
\workings 到达相差 $-2\pi(4.5-3)=-3\pi$，等价于π，等幅反相合幅0。
\crosscheck 相差为奇数半波长，且原题明确振幅不变。
\sourceproblem{填空19}{185}{191}
\problemtext 两等幅A相干源S1、S2相隔半波长，S1超前S2 $\pi/2$；图中P在S1外侧，求合幅。
\workings $r_2-r_1=\lambda/2$，到达相差 $\phi_{1P}-\phi_{2P}=\pi/2+k(r_2-r_1)=3\pi/2$。若到达幅值均为A，则 $A_P=\sqrt{2A^2+2A^2\cos(3\pi/2)}=\sqrt2 A$。
\crosscheck 若传播幅值随路程变化，需用到达幅值 $a_1,a_2$，结果为 $\sqrt{a_1^2+a_2^2}$；不能由源幅相等推出任意波的到达幅相等。
\sourceproblem{填空20}{185}{191}
\problemtext 两右行相干波源A、B相距 $3\lambda/2$，A在正峰时B在负峰，P在B右侧，求合幅。
\workings 源相差π，传播路程差 $r_A-r_B=3\lambda/2$，故到达相差 $\pi-3\pi=-2\pi$，同相。按无衰减平面波模型，合幅 $A_1+A_2$。
\crosscheck 无吸收本身不排除几何扩散；若有扩散，应相加实际到达幅值。
\sourceproblem{填空21}{185}{191}
\problemtext 驻波 $y=1.0\cos(\pi x/3)\cos40\pi t$，x、y用cm，求分波速度及相邻节点间距。
\workings $\omega=40\pi$ rad/s，$k=\pi/3$ rad/cm，$u=120$ cm/s，$\lambda=6$ cm，相邻节点距3 cm。
\crosscheck 节点为 $x=1.5+3j$ cm；波速不是质元振动速率。
\sourceproblem{填空22}{185}{191}
\problemtext 入射波 $A\cos[2\pi(\nu t+x/\lambda)]$ 在x=0反射，写固定端及自由端反射波。
\workings 理想完全反射固定端：$y_r=-A\cos[2\pi(\nu t-x/\lambda)]$；自由端：$y_r=A\cos[2\pi(\nu t-x/\lambda)]$。
\crosscheck 固定端总位移恒零；自由端总空间导数恒零，分别满足理想弦边界条件。
\sourceproblem{填空23}{185}{191}
\problemtext 驻波 $y=0.1\cos(\pi x)\cos(90\pi t)$（SI），求分波波长及频率。
\workings $k=\pi$ rad/m，$\omega=90\pi$ rad/s，故 $\lambda=2$ m、$\nu=45$ Hz。
\crosscheck 两分波各振幅0.05 m；节点间距1 m。
\sourceproblem{填空24}{185}{191}
\problemtext 驻波 $2A\cos(2\pi x/\lambda)\cos(2\pi\nu t)$ 的相邻节点间距。
\workings 节点为 $x=\lambda/4+j\lambda/2$，相邻间距 $\lambda/2$。
\crosscheck 节点至最近腹点的距离为 $\lambda/4$，并非节点间距。
\sourceproblem{填空25}{185}{191}
\problemtext 驻波 $y=A\cos2\pi x\cos100\pi t$（SI），求x=1/8、3/8 m处相差。
\workings 空间因子分别为 $\cos(\pi/4)=\sqrt2/2$ 和 $\cos(3\pi/4)=-\sqrt2/2$，故两点反相，差π。
\crosscheck 驻波相位不按行波的 $k\Delta x$ 直接计算；两点之间有一节点。
\sourceproblem{填空26}{185}{191}
\problemtext 入射波 $A\cos[2\pi(t/T+x/\lambda)]$ 在x=0自由端等强反射，求反射式和 $x=\lambda/3$ 处合幅。
\workings 反射波 $y_r=A\cos[2\pi(t/T-x/\lambda)]$。总波 $2A\cos(2\pi x/\lambda)\cos(2\pi t/T)$，故合幅 $|2A\cos(2\pi/3)|=A$。
\crosscheck 负空间因子意味着相位反转，合幅必须非负；自由端总斜率为零。
\sourceproblem{填空27}{185}{191}
\problemtext 驻波同一节点两侧等距的非节点质元相位差。
\workings 理想单频驻波的空间因子在节点两侧等距处互为相反数，故相位差π（模 $2\pi$）。
\crosscheck 若两个点恰又落在其他节点，则振幅为零，相位无定义。
\sourceproblem{填空28}{185}{191}
\problemtext 50 Hz波、波速350 m/s，两点相位差 $2\pi/3$，求间距。
\workings $\lambda=7$ m，取未折返相差大小或最小正间距，$d=\lambda/3=7/3$ m。
\crosscheck 若仅指定模 $2\pi$ 的无向相位差，则还允许 $|j\lambda\pm\lambda/3|$ 的非零间距。
\sourceproblem{填空29}{185}{191}
\problemtext 火车速20 m/s，汽笛600 Hz，声速340 m/s，求静止观察者在车前及车后听到的频率。
\workings 移动声源公式 $f_{\rm 前}=600(340)/(340-20)=637.5$ Hz，$f_{\rm 后}=600(340)/(340+20)\approx566.67$ Hz。
\crosscheck 假设静止空气、沿径向传播，前方趋近频率较高。
\sourceproblem{填空30}{185}{191}
\problemtext 静止警报器1000 Hz，汽车以79.2 km/h趋近或远离，声速340 m/s，求车内频率。
\workings 观察者速度22 m/s，$f_{\rm 近}=1000(340+22)/340\approx1064.71$ Hz，$f_{\rm 远}=1000(340-22)/340\approx935.29$ Hz。
\crosscheck 移动观察者速度在分子，与移动声源公式有别。
\sourceproblem{计算1}{186}{192}
\problemtext 由平面波 $y=\cos[\pi(100t-2x)+3\pi/2]$ 求振幅、波速、周期、波长及传播方向。
\workings 按SI理解，$A=1$ m，$\omega=100\pi$ rad/s，$k=2\pi$ rad/m，故 $u=50$ m/s、$T=0.02$ s、$\lambda=1$ m，沿正x方向。
\crosscheck 原题未标明单位体系，所列数值按SI；固定相位给正传播速度。
\sourceproblem{计算2}{186}{192}
\problemtext 波源 $y_0=4\times10^{-3}\cos240\pi t$ m，波速30 m/s，求周期、波长及波函数。
\workings $T=1/120$ s，$\lambda=30/120=0.25$ m。选传播方向为正x、源为原点，$y=4\times10^{-3}\cos(240\pi t-8\pi x)$ m。
\crosscheck 若坐标正向取反则空间相位符号改变；$x=0$ 时恢复源振动。
\sourceproblem{计算3}{186}{192}
\problemtext 绳上波 $y=20\cos[\pi(2.50t-0.01x)]$，x、y单位cm，求波参量、最大质元速度、1 s及2 s波形、0.2 m和0.7 m处相差。
\workings $A=20$ cm，$\omega=2.5\pi$ rad/s，$k=0.01\pi$ rad/cm，$u=250$ cm/s，$\nu=1.25$ Hz，$\lambda=200$ cm；最大质元速率 $A\omega=50\pi$ cm/s。两时刻波形分别为
\[
y(x,1)=20\sin(\pi x/100),\qquad y(x,2)=-20\cos(\pi x/100).
\]
绘图取x=0,50,100,150,200 cm，第一曲线纵坐标依次为0,20,0,-20,0 cm，第二为-20,0,20,0,-20 cm；按正弦周期延伸。两点间距50 cm，后点相对前点的未折返相差为 $-\pi/2$。
\crosscheck 波形1秒后沿正x平移250 cm，模一波长为50 cm；不得将cm波数直接与m坐标相乘。
\sourceproblem{计算4}{186}{192}
\problemtext 波源周期0.02 s、波速100 m/s，初时正向过零，求距源15、5 m质元振动及初相，另求16、17 m处相差。
\workings 设题中未给振幅为A，源 $y_0=A\sin100\pi t$，$k=\pi$ rad/m。选正传播方向，$y=A\sin(100\pi t-\pi x)$。15、5 m处均为 $-A\sin100\pi t=A\cos(100\pi t+\pi/2)$，余弦初相均为 $\pi/2$ 模 $2\pi$。17 m处相对16 m处相差 $-\pi$。
\crosscheck 两点15、5 m相隔5个波长而同相；若振动从t=0才启动，公式仅在各点波到达后成立，初相是稳态外推约定。
\sourceproblem{计算5}{186}{192}
\problemtext 左行波 $y=2\times10^{-3}\cos(10t+6x)$（SI），求峰谷经过原点的时刻及t=6 s时各峰谷坐标。
\workings 原点正峰满足 $10t=2j\pi$，即 $t=j\pi/5$；负峰 $t=(2j+1)\pi/10$。若只考虑非负时间，取使t非负的整数j。t=6 s时正峰 $x=(2j\pi-60)/6$ m，负峰 $x=[(2j+1)\pi-60]/6$ m。
\crosscheck 相邻同类峰距 $\pi/3$ m，峰谷相隔 $\pi/6$ m，固定相位向负x传播。
\sourceproblem{计算6}{187}{193}
\problemtext 右行绳波 $y=0.05\cos(40\pi t-4\pi x)$（SI），求周期、波速、距源0.1 m处振动及0.2 s波形。
\workings $T=0.05$ s，$u=10$ m/s，$\lambda=0.5$ m；$y(0.1,t)=0.05\cos(40\pi t-0.4\pi)$ m。0.2 s为4个周期，波形 $y=0.05\cos4\pi x$ m。
\crosscheck 若波源在t=0才启动，受扰区限于 $0\le x\le2$ m；题给稳态波式未明说启动边界。
\sourceproblem{计算7}{187}{193}
\problemtext 右行波，速度4 m/s、周期0.01 s、振幅0.01 m，初始图原点位移A/2且空间斜率负，求波函数及原点振动。
\workings $\omega=200\pi$ rad/s，$\lambda=0.04$ m，$k=50\pi$ rad/m。原点初速度为正，故 $\phi=-\pi/3$，波函数 $y=0.01\cos(200\pi t-50\pi x-\pi/3)$ m，原点式 $y_0=0.01\cos(200\pi t-\pi/3)$ m。
\crosscheck 初位移0.005 m；初时空间斜率负、时间导数正，符合原图及右行关系。
\sourceproblem{计算8}{187}{193}
\problemtext $y=0.25\cos(125t-0.37x)$（SI），求x=10、25 m的振动式、相差和x=10 m在4 s的位移。
\workings $y_{10}=0.25\cos(125t-3.7)$ m，$y_{25}=0.25\cos(125t-9.25)$ m。后点相对前点的未折返相差为 $-5.55$ rad，主值为 $2\pi-5.55$ rad。$y(10,4)=0.25\cos496.3$ m，约 $0.249359$ m。
\crosscheck 相位角以弧度计；未折返相差与主值不可混同。
\sourceproblem{计算9}{187}{193}
\problemtext 初始波形图中原点位移 $\sqrt2 A/2$ 且斜率负，P在下降段且向下运动，相邻上升、下降过零点相距100 m，频率250 Hz，求波函数、100 m处振动与速度及四分之一周期波形。
\workings P处速度与空间斜率同为负，波向负x传播；$\lambda=200$ m，$k=\pi/100$ rad/m，$\omega=500\pi$ rad/s，$u=5.0\times10^4$ m/s。原点初相 $\pi/4$，故
\[
y=A\cos(500\pi t+\pi x/100+\pi/4),\quad
y(100,t)=A\cos(500\pi t+5\pi/4),
\]
\[
v(100,t)=-500\pi A\sin(500\pi t+5\pi/4).
\]
在 $t=T/4=0.001$ s，波形 $y=A\cos(\pi x/100+3\pi/4)$，相对原图左移50 m。绘图关键点：x=-75、-25、25、75、125 m处位移依次为A、0、-A、0、A。
\crosscheck 原题未给A数值，保留参数；原点初位移为正且空间斜率负。波速是按原图标尺和给定频率所得形式值，未假设具体介质。

\sourceproblem{计算10}{188}{194}
\problemtext 图示右行波初始振幅0.02 m，原点上升过零，x=2 m下降过零，波速10 m/s，求原点、全波和x=3 m处振动式。
\workings 图示波长4 m，$k=\pi/2$ rad/m，$\omega=uk=5\pi$ rad/s。原点初速度为负，故
\[
y_O=-0.02\sin5\pi t,\quad
y(x,t)=0.02\sin(\pi x/2-5\pi t),\quad
y_P=-0.02\cos5\pi t
\]
（SI）。
\crosscheck 初始波形在x=1 m为正峰、x=3 m为负峰；$\partial_t y=-u\partial_x y$，传播方向正确。
\sourceproblem{计算11}{188}{194}
\problemtext 已知右行波 $y=A\cos[2\pi(\nu t-x/\lambda)+\phi]$，图示P在原点左方L，求其振动、速度和加速度。
\workings 记 $\Theta=2\pi(\nu t+L/\lambda)+\phi$，则
\[
y_P=A\cos\Theta,\quad v_P=-2\pi\nu A\sin\Theta,\quad
a_P=-(2\pi\nu)^2A\cos\Theta.
\]
\crosscheck P坐标为负L，不能以正距离直接代入x；加速度始终满足简谐运动关系。
\sourceproblem{计算12}{188}{194}
\problemtext 波源周期0.01 s，初时正向过零，波速400 m/s，求距源8 m处振动和初相，以及9、10 m处相差。
\workings 设振幅为题中未给的A，$\omega=200\pi$ rad/s，$\lambda=4$ m。沿传播方向取正坐标，$y=A\sin(200\pi t-\pi x/2)$。故 $y_8=A\sin200\pi t$，余弦式初相 $-\pi/2$（模 $2\pi$）。后一点相对前一点未折返相差为 $-\pi/2$。
\crosscheck 8 m为两个波长；若波源在零时刻才启动，稳态式仅在波到达后有效。原题未给振幅，不指定数值。
\sourceproblem{计算13}{188}{194}
\problemtext 右行横波波速100 m/s，每米含54个波长，振幅0.03 m，原点初时正向过零，求波函数。
\workings $\lambda=1/54$ m，$\nu=5400$ Hz，$k=108\pi$ rad/m，$\omega=10800\pi$ rad/s。故
\[
y=0.03\sin(10800\pi t-108\pi x)\quad(\mathrm m).
\]
\crosscheck 原点初速度为正；$\omega/k=100$ m/s。54是空间周期数每米，并非波长54 m。
\sourceproblem{计算14}{189}{195}
\problemtext 300 Hz平面声波在截面积0.0300平方米的空气管中以340 m/s传播，10 s通过能量0.0270 J，求平均能流、能流密度和平均能量密度。
\workings $\overline P=\Delta E/\Delta t=2.70\times10^{-3}$ W；$\overline I=\overline P/S=0.0900$ W每平方米；对无耗散单向平面行波 $\overline w=\overline I/u=2.65\times10^{-4}$ J每立方米。
\crosscheck 三者分别为功率、面积功率密度、体积能量密度；频率不参与本题能量换算。假定截面强度均匀，否则为截面平均值。
\sourceproblem{计算15}{189}{195}
\problemtext 反相点波源A、B相距30 cm，PB为40 cm且垂直AB；求P处干涉最大限度削弱时最长波长。
\workings $AP=\sqrt{30^2+40^2}=50$ cm，路程差10 cm。源相差π，到达反相条件为 $2\pi(10\,\mathrm{cm})/\lambda=2n\pi$，其中正整数n。因此 $\lambda=10/n$ cm，最长10 cm。
\crosscheck 最大限度削弱指相位相反，未必完全抵消；点源传播振幅若不相等，残余幅值为两到达幅值差的绝对值。
\sourceproblem{计算16}{189}{195}
\problemtext 图示同相波源P、Q相隔 $3\lambda/2$，R在Q外侧，两源同时发出等幅A的相干波，求R处相差与合幅。
\workings 按图 $PR-QR=3\lambda/2$，故P波相对Q波的未折返相差为 $-3\pi$，等价于π。若两分波到R的幅值均为A，则 $A_R=0$。
\crosscheck 原图R位于两源连线延长线上；若R在两源之间，路程差随位置改变。若源幅相等但传播幅值不同，合幅应为 $|a_P-a_Q|$，不一定为零。
\sourceproblem{计算17}{189}{195}
\problemtext 第二平面波源超前第一 $\pi/4$，波长8 m，至P路程12、14 m，到达幅值0.30、0.20 m，求合幅。
\workings 到达相差 $\Delta\phi=\pi/4-(2\pi/8)(14-12)=-\pi/4$。故
\[
A_P=\sqrt{0.30^2+0.20^2+2(0.30)(0.20)\cos(\pi/4)}
=\sqrt{0.13+0.06\sqrt2}\ \mathrm m
\approx0.4635\ \mathrm m .
\]
\crosscheck 合幅在0.10至0.50 m之间；题给的是到达幅值，无需另作距离衰减修正。
\sourceproblem{计算18}{190}{196}
\problemtext 两反相波源间距30 cm，频率10 Hz，PB垂直AB且长40 cm，求P处干涉极大对应的波长。
\workings 路程差 $AP-BP=50-40=10$ cm。到达同相要求 $\pi-2\pi(10\,\mathrm{cm})/\lambda=2j\pi$，等价于
\[
\lambda=\frac{20}{2n+1}\ \mathrm{cm},\qquad n=0,1,2,\ldots .
\]
最大允许波长20 cm；相应波速为 $u=2/(2n+1)$ m/s。
\crosscheck 频率已知但波速未给，不能从题设选定唯一阶次；若振幅随传播距离变化，同相仍给固定到达幅值下的最大合幅。
\sourceproblem{计算19}{190}{196}
\problemtext 两相干反相源相距30 m，频率100 Hz、波速400 m/s，连线方向两波等幅不衰减，求两源之间静止点。
\workings $\lambda=4$ m，以S1为x=0、S2为x=30 m，$r_1=x,r_2=30-x$。反相抵消条件为 $r_2-r_1=4n$ m（整数n），故 $x=15-2n$ m。限制 $0<x<30$ 得 $n=-7,\ldots,7$，即 $x=1,3,5,\ldots,29$ m，共15点。
\crosscheck 中点15 m等程且源反相，必为节点；节点间距2 m，端点不计入。
\sourceproblem{计算20}{190}{196}
\problemtext 同绳两波 $0.06\cos[\pi(x-4t)]$、$0.06\cos[\pi(x+4t)]$（SI），求分波参数、驻波节点波腹及指定位置振幅。
\workings $k=\pi$ rad/m、$\omega=4\pi$ rad/s，故 $\lambda=2$ m、$\nu=2$ Hz、$u=4$ m/s，第一波右行、第二波左行。叠加得
\[
y=0.12\cos\pi x\cos4\pi t .
\]
节点 $x=n+1/2$ m，波腹 $x=n$ m（整数n）。腹点幅0.12 m，x=1.2 m处幅 $0.12|\cos1.2\pi|\approx0.09708$ m。
\crosscheck 驻波空间因子为负时相位反转而非负振幅；相邻节点距1 m。
\sourceproblem{计算21}{190}{196}
\problemtext 给定右行波 $0.05\cos[2\pi(t/0.05-x/4)]$，等幅反向波在原点同相，求反向波、合驻波及四个最近节点。
\workings 反向波 $y_2=0.05\cos(40\pi t+\pi x/2)$，合成 $y=0.10\cos(\pi x/2)\cos40\pi t$（SI）。节点满足 $x=1+2n$ m，最近四点为 $-3,-1,1,3$ m。
\crosscheck 若绳实际仅定义于正半轴，则最近四个可取1、3、5、7 m；题图未给有限绳边界，以全轴模型为主。
\sourceproblem{计算22}{191}{197}
\problemtext O点波源向两侧发出等幅A、频率 $\nu$ 的波，左侧波密界面距O为 $5\lambda/4$，全反射，求O两侧合成波函数。
\workings 取O为原点、向右为正，界面 $x=-d$，令 $\Theta=\omega t+\phi$、$k=2\pi/\lambda$。源左行波 $A\cos(\Theta+kx)$，右行波 $A\cos(\Theta-kx)$。按教材波密界面理想位移反相全反射模型，反射波为 $-A\cos(\Theta-kx-2kd)$。因 $2kd=5\pi$，反射波化为 $A\cos(\Theta-kx)$。于是
\[
y(x,t)=
\begin{cases}
2A\cos kx\cos\Theta,&-d\le x\le0,\\
2A\cos(\Theta-kx),&x\ge0 .
\end{cases}
\]
\crosscheck 界面处 $\cos kd=\cos(5\pi/2)=0$，总位移为零，原点两侧连续；假定反射波透过理想源位置且源发射幅相固定。若实际界面反射系数并非-1或波源散射反射波，则需重新指定边界条件。
\sourceproblem{计算23}{191}{197}
\problemtext 右行波在固定点B反射，A坐标 $\lambda/2$，AB长 $\lambda/2$，A处入射振动为 $A\cos\omega t$，求入射、反射及驻波。
\workings 记 $k=2\pi/\lambda$，B坐标为 $\lambda$，用幅值A表示原题振幅。入射波 $y_i=A\cos(\omega t-kx+\pi)$。固定端要求总位移为零，故反射波 $y_r=A\cos(\omega t+kx)$。合成
\[
y=-2A\sin\omega t\sin kx .
\]
\crosscheck A处入射相位为 $\omega t$；B处两分波相反且总位移恒零，反射波向负x传播。
\sourceproblem{计算24}{191}{197}
\problemtext 1 kHz声源以34 m/s向右，右方反射面以68 m/s向左，声速340 m/s，求前方入射波长、每秒到达反射面波数和反射波长。
\workings 静止空气中入射波长 $\lambda_i=(340-34)/1000=0.306$ m；反射面与入射波峰相向，接收频率 $f_m=(340+68)/0.306=1333.33$ Hz。反射面作为向左运动的相干次级波源，左行反射波长
\[
\lambda_r=(340-68)/f_m=0.204\ \mathrm m .
\]
相应空气中反射波频率为 $340/0.204\approx1666.67$ Hz。
\crosscheck 反射面参考系中的入射与反射频率相同；不能把其接收频率直接当作空气系反射频率。仅适用于声源与反射面尚未相遇、均匀静止空气和理想镜面反射模型。
\sourceproblem{应用1}{191}{197}
\problemtext 固定超声源频率100 kHz，迎面驶来汽车反射回波110 kHz，声速330 m/s，求车速。
\workings 车作为移动接收者，接收频率 $f_1=f_0(u+v)/u$；作为向源运动的反射次级波源，回波频率 $f_r=f_1u/(u-v)=f_0(u+v)/(u-v)$。解得
\[
v=u\frac{f_r-f_0}{f_r+f_0}
=330\frac{10}{210}=15.7143\ \mathrm{m/s}
\]
约56.57 km/h。
\crosscheck 代回频率比为1.10；拍频为10 kHz而不是回波频率110 kHz。假设径向运动、静止空气及相干反射。

\sourceproblem{应用2}{192}{198}
\problemtext 干涉消声器将300 Hz声波分为直管和弯管两路，在另一端汇合，声速340 m/s，求最小管长差。
\workings 假定分流与汇合不引入额外相差，反相条件为 $\Delta r=(2n+1)\lambda/2$，$\lambda=u/\nu=340/300$ m，最小正值为 $17/30$ m，约0.567 m。
\crosscheck 完全相消还要求两路到达幅值相等；管道损耗、端部相移及宽频噪声会限制实际消声效果。
\sourceproblem{应用3}{192}{198}
\problemtext 警笛800 Hz，车速25 m/s，声速330 m/s；求静止观察者在趋近及远离时听到的频率，以及前方同向15 m/s观察者的频率。
\workings 静止观察者听到 $f_+=800(330)/(330-25)\approx865.57$ Hz，$f_-=800(330)/(330+25)\approx743.66$ Hz。同向前方观察者远离声源发出的前行波峰，$f'=800(330-15)/(330-25)\approx826.23$ Hz。
\crosscheck 源与观察者速度分别相对静止空气取值，不能简单把相对速度10 m/s代入移动声源公式；假定沿传播直线运动。
\sourceproblem{应用4}{192}{198}
\problemtext 水中两个观察平台沿同一直线相向运动，速度50.0和70.0 km/h，第一平台发声频率1000 Hz，声速5470 km/h，求第二平台的接收频率和第一平台收到的反射频率。
\workings 统一采用km/h，正向声波的移动源和移动观察者频移给 $f_2=1000(5470+70)/(5470-50)\approx1022.14$ Hz。反射波在第二平台系与入射同频；将该平台视为向第一平台运动的次级声源，回波频率 $f_{\rm echo}=f_2(5470+50)/(5470-70)\approx1044.85$ Hz。
\crosscheck 两次频移均为趋近，所得频率逐次升高；使用静止均匀水介质、经典径向运动及理想相干反射模型。
\sourceproblem{应用5}{192}{198}
\problemtext 窗口面积1.0平方米，声强级80 dB，求经窗口进入的声功率。
\workings 取标准参考声强 $I_0=10^{-12}$ W每平方米，$I=I_0\,10^{80/10}=10^{-4}$ W每平方米。若声强均匀且垂直入射、窗口无损透射，则 $P=IS=10^{-4}$ W。
\crosscheck 声强级为对数值而非功率；若声波斜入射或窗口存在反射，应积分法向声强并计透射率，题给数据仅确定理想模型下的值。
